\section{Study 2: Arm-Side Stochasticity}
\label{sec:study2}

% ─────────────────────────────────────────────────────────────────────────────
\subsection{Motivation}
\label{sec:study2:motivation}

The NumPy ballistic simulator is deterministic: given a commanded velocity
$\mathbf{v}_{\mathrm{cmd}}$, the landing position is a fixed function of that input.
A closed-form ballistic inverter could therefore solve the task in one step,
rendering the GP-based learning loop superfluous.
In a physical system, the gap between commanded and delivered velocity arises from
several sources: joint-torque saturation, gripper compliance, and release timing
uncertainty.
\mcpilot{}~\cite{turcato2025mcpilot} explicitly models the last source using a
gripper delay parameter estimated by Bayesian optimisation (\S5 of the original paper).
We replace that calibration step with an explicit stochastic noise model injected at
the velocity-override interface (Section~\ref{sec:system}), which lets us study the
effect of arm-side stochasticity in a controlled, parameter-varied manner while
preserving exact GP-gradient compatibility.

% ─────────────────────────────────────────────────────────────────────────────
\subsection{Three Noise Models}
\label{sec:study2:models}

All three models are realised through the \texttt{perturb\_numpy} interface:
$\texttt{perturb\_numpy}(\mathbf{v}_{\mathrm{cmd}},\, n) \to (\alpha,\, \boldsymbol{\delta})$,
returning a scalar scale $\alpha$ and an additive offset $\boldsymbol{\delta} \in \mathbb{R}^3$
such that

\begin{equation}
  \mathbf{v}_{\mathrm{actual}} = \alpha\,\mathbf{v}_{\mathrm{cmd}} + \boldsymbol{\delta}.
  \label{eq:perturb_interface}
\end{equation}

The scale $\alpha$ and offset $\boldsymbol{\delta}$ are \emph{detached} constants
(no gradient tape), so the policy gradient flows through $\mathbf{v}_{\mathrm{cmd}}$
unobstructed — a necessary condition for correct Monte Carlo policy-gradient
estimation.

\paragraph{VelocityBiasNoise.}
A zero-mean additive perturbation:
\begin{equation}
  \alpha = 1, \quad \boldsymbol{\delta} \sim \mathcal{N}(\mathbf{0},\, \sigma^2 I_3).
  \label{eq:bias_noise}
\end{equation}
This models isotropic sensor noise or small random errors in the velocity-override
step.  The noise has zero mean, so the expected delivered velocity equals the commanded
velocity.

\paragraph{VelocitySlipNoise.}
A multiplicative under-delivery with additive scatter:
\begin{equation}
  \alpha = 1 - \alpha_{\mathrm{slip}}, \quad \boldsymbol{\delta} \sim \mathcal{N}(\mathbf{0},\, \sigma^2 I_3),
  \label{eq:slip_noise}
\end{equation}
where $\alpha_{\mathrm{slip}} \in (0, 1)$ is the slip fraction.  The expected
delivered speed is $(1 - \alpha_{\mathrm{slip}})\|\mathbf{v}_{\mathrm{cmd}}\|$,
introducing a systematic undershoot that a sufficiently informed policy can compensate.

\paragraph{ReleaseTimingJitter (paper-faithful).}
A gripper delay $t_d \sim \mathcal{U}(a, b)$ with $a > 0$ causes the arm to decelerate
for $t_d$ seconds beyond the commanded release instant.
The delivered velocity is approximately
\begin{equation}
  \mathbf{v}_{\mathrm{actual}} \approx
    \underbrace{\left(1 - \frac{\gamma\,t_d}{\|\mathbf{v}_{\mathrm{cmd}}\|}\right)}_{\alpha}
    \mathbf{v}_{\mathrm{cmd}},
  \label{eq:jitter_noise}
\end{equation}
where $\gamma$ is the arm deceleration rate.  This matches the gripper-delay model of
\cite{turcato2025mcpilot} and sets $\boldsymbol{\delta} = \mathbf{0}$.

% ─────────────────────────────────────────────────────────────────────────────
\subsection{Zero-Mean Unlearnability and GP Degeneration}
\label{sec:study2:theory}

\begin{proposition}[Optimal Action Invariance Under Symmetric Noise]
\label{prop:zeromean}
Let $f: \mathbb{R}^3 \to \mathbb{R}^2$ be the deterministic ballistic map from release
velocity to landing position, and let $\boldsymbol{\epsilon} \sim p_\epsilon$ with
$\mathbb{E}[\boldsymbol{\epsilon}] = \mathbf{0}$.
Define the noise-aware objective
$J_{\mathrm{aware}}(\theta) = \mathbb{E}_{\boldsymbol{\epsilon}}[c(f(\mathbf{v}_{\mathrm{cmd}}
+ \boldsymbol{\epsilon}), \mathbf{P})]$
and the noiseless objective
$J_{\mathrm{naive}}(\theta) = c(f(\mathbf{v}_{\mathrm{cmd}}), \mathbf{P})$.
If $p_\epsilon$ is symmetric around zero, the minimisers of $J_{\mathrm{aware}}$ and
$J_{\mathrm{naive}}$ coincide.
That is, symmetric zero-mean noise does not shift the optimal mean action of the policy,
even though the GP correctly learns the increased aleatoric uncertainty it induces.
\end{proposition}

\begin{proof}[Proof sketch]
The cost \eqref{eq:cost} is a function of $\|\mathbf{p}_T - \mathbf{P}\|^2$ with
minimum at $\mathbf{p}_T = \mathbf{P}$.
For symmetric $p_\epsilon$, $\mathbb{E}[f(\mathbf{v}^* + \boldsymbol{\epsilon})]
= f(\mathbf{v}^*)$ whenever $f$ is locally affine near $\mathbf{v}^*$
(valid to first order by Taylor expansion).
The gradient condition
$\nabla_{\mathbf{v}_{\mathrm{cmd}}} J_{\mathrm{aware}} = 0$ therefore reduces to
$\mathbb{E}_{\boldsymbol{\epsilon}}[\nabla_{\mathbf{v}} c \big|_{\mathbf{v}_{\mathrm{cmd}}
+ \boldsymbol{\epsilon}}] = 0$, which is satisfied at the same $\mathbf{v}^*$ as the
noiseless problem due to the symmetry of $p_\epsilon$.
\end{proof}

The practical implication is that a \emph{naive} policy (optimised assuming
$\boldsymbol{\epsilon} = \mathbf{0}$) and a \emph{noise-aware} policy reach the same
optimum when the noise is zero-mean.
The difference lies only in cost calibration: the naive policy evaluates 400 Monte
Carlo particles without noise and therefore overestimates its own precision
(falsely low cost value), whereas the aware policy propagates noise through the
MC rollouts and reports an honest expected cost.

\paragraph{GP degeneration under exact data.}
An unexpected consequence of zero-mean noise appears in GP maximum likelihood
estimation.
With $\sigma = 0$ (noise completely removed), $N_{\mathrm{exp}} = 10$ perfectly
deterministic exploration throws give the GP log-likelihood a degenerate ridge: the
MLE optimizer drives lengthscales to 1791\,m, 1602\,m, and 357\,m — three to four
orders of magnitude larger than the 1\,m-scale domain.
A GP with these lengthscales is equivalent to a constant function: it cannot
distinguish throws at different commanded speeds.
With $\sigma = 0.04$, repeated throws at the same speed produce slightly different
landing positions, so the MLE finds finite lengthscales that explain the
position variance as a function of speed.
This reveals a counterintuitive role for noise: a small amount of stochasticity is
\emph{necessary} for GP identifiability when the exploration data are not
independently varied enough to avoid the constant-function degenerate solution.

% ─────────────────────────────────────────────────────────────────────────────
\subsection{Empirical Comparison: Noise-Aware vs.\ Naive Policy}
\label{sec:study2:empirical}

We compare noise-aware and naive training under VelocitySlipNoise with
$\alpha_{\mathrm{slip}} = 0.20$, $\sigma = 0.04$, $u_M = 3.0$\,m/s, seed=1, and 10
evaluation trials (throws 6–15, after 5 shared exploration throws).

\begin{table}[t]
  \centering
  \caption{Aggregate performance of noise-aware and naive policies under 20\%
           velocity-slip noise ($\sigma = 0.04$).
           ``Total'' counts all 15 throws (5 exploration + 10 evaluation);
           ``policy-only'' counts evaluation throws 6–15 only.}
  \label{tab:noise_aggregate}
  \small
  \begin{tabular}{lcccc}
    \toprule
    Policy & Total hits & Policy-only hits & Mean error & Final cost \\
    \midrule
    Aware & 8/15 (53\%) & 5/10 (50\%) & 0.100\,m & $\approx 0.009$ \\
    Naive & 4/15 (27\%) & 1/10 (10\%) & 0.125\,m & $\approx 0.001$ \\
    \bottomrule
  \end{tabular}
\end{table}

\begin{table}[t]
  \centering
  \caption{Per-throw landing errors for evaluation throws 6–15 (in metres).
           Checkmarks denote hits (error $< 0.10$\,m). Naive cost $\approx 0.001$
           is falsely optimistic: it is evaluated without noise.}
  \label{tab:noise_pershot}
  \small
  \begin{tabular}{ccc}
    \toprule
    Throw & Aware & Naive \\
    \midrule
    6  & 0.094\,$\checkmark$ & 0.129\,$\times$ \\
    7  & 0.039\,$\checkmark$ & 0.151\,$\times$ \\
    8  & 0.124\,$\times$     & 0.135\,$\times$ \\
    9  & 0.084\,$\checkmark$ & 0.153\,$\times$ \\
    10 & 0.106\,$\times$     & 0.090\,$\checkmark$ \\
    11 & 0.103\,$\times$     & 0.107\,$\times$ \\
    12 & 0.033\,$\checkmark$ & 0.106\,$\times$ \\
    13 & 0.106\,$\times$     & 0.101\,$\times$ \\
    14 & 0.118\,$\times$     & 0.141\,$\times$ \\
    15 & 0.055\,$\checkmark$ & 0.127\,$\times$ \\
    \midrule
    Hits & 5/10 & 1/10 \\
    \bottomrule
  \end{tabular}
\end{table}

\paragraph{Interpretation.}
The aware policy achieves 50\% hits versus 10\% for the naive policy — a factor of
five improvement.
The mechanism is speed compensation: the systematic 20\% undershoot means that to
land at distance $d$ the arm must command speed $v^* / (1 - \alpha_{\mathrm{slip}}) =
1.25\,v^*$.
The noise-aware policy learns to output speeds $\approx 1.25\times$ larger than the
noiseless optimum because its MC rollouts propagate the slip factor; the naive policy,
optimising without noise, targets $v^*$ and consistently undershoots.

The 50\% ceiling on aware performance is not a learning failure — it reflects
irreducible scatter from the $\sigma = 0.04$ isotropic additive component.
At a release speed of $\approx 2.5$\,m/s, $\sigma = 0.04$\,m/s translates to a
landing position standard deviation of roughly 2–4\,cm, which alone accounts for
half the throws missing the 10\,cm-radius target.
Multi-seed statistics would be needed to confirm this bound; we leave that
to \needdata{multi-seed noise comparison ($\geq 5$ seeds per condition)}.

The naive policy's final cost of $\approx 0.001$ versus the aware policy's
$\approx 0.009$ is a direct manifestation of the false-optimism consequence of
Proposition~\ref{prop:zeromean}: the naive optimiser sees noiseless rollouts and
concludes it has solved the task precisely, while the aware policy honestly reports
the expected cost under the noise distribution it was trained with.
This discrepancy — lower training cost yet far worse test performance — is
a concrete demonstration of the distributional mismatch introduced by ignoring
physical uncertainty during policy optimisation.
